
\documentclass{article}
\usepackage{colortbl}
\usepackage{makecell}
\usepackage{multirow}
\usepackage{supertabular}

\begin{document}

\newcounter{utterance}

\centering \large Interaction Transcript for game `cladder', experiment `full\_v1.5\_default', episode 4570 with Qwen3.6{-}27B.
\vspace{24pt}

{ \footnotesize  \setcounter{utterance}{1}
\setlength{\tabcolsep}{0pt}
\begin{supertabular}{c@{$\;$}|p{.15\linewidth}@{}p{.15\linewidth}p{.15\linewidth}p{.15\linewidth}p{.15\linewidth}p{.15\linewidth}}
    \# & \multicolumn{2}{c}{Player} && \multicolumn{2}{c}{Game Master} \\
    \hline

    \theutterance \stepcounter{utterance}  
    & & & \multicolumn{4}{p{0.6\linewidth}}{
        \cellcolor[rgb]{0.9,0.9,0.9}{
            \makecell[{{p{\linewidth}}}]{
                \texttt{\tiny{[P1$\langle$GM]}}
                \texttt{You are an expert in causal inference. The following question is not a typical commonsense query, but rather a meticulously designed question created by a professor specializing in causal inference, intended to assess the students' mastery of the course content.} \\
\texttt{Imagine a self{-}contained, hypothetical world with only the following conditions, and without any unmentioned factors or causal relationships: Poverty has a direct effect on liking spicy food and cholera. Water company has a direct effect on liking spicy food. Liking spicy food has a direct effect on cholera. Poverty is unobserved. The overall probability of liking spicy food is 41\%. For people who do not like spicy food, the probability of cholera contraction is 59\%. For people who like spicy food, the probability of cholera contraction is 33\%. Is cholera contraction more likely than cholera prevention overall?} \\
\texttt{Start your answer with 'Yes' or 'No', followed by additional reasoning or evidence to support your explanation.} \\
            }
        }
    }
     \\ \\

    \theutterance \stepcounter{utterance}  
    & \multicolumn{4}{p{0.6\linewidth}}{
        \cellcolor[rgb]{0.561,0.737,0.561}{
            \makecell[{{p{\linewidth}}}]{
                \texttt{\tiny{[P1$\rangle$GM]}}
                \texttt{Yes} \\
\\ 
\texttt{To determine whether cholera contraction is more likely than cholera prevention overall, we need to calculate the total probability of contracting cholera, denoted as \$P(C)\$.} \\
\\ 
\texttt{We are given the following information:} \\
\texttt{1.  The overall probability of liking spicy food: \$P(S) = 0.41\$.} \\
\texttt{2.  Consequently, the probability of not liking spicy food is: \$P(\textbackslash{}neg S) = 1 {-} P(S) = 1 {-} 0.41 = 0.59\$.} \\
\texttt{3.  The probability of cholera given no liking of spicy food: \$P(C | \textbackslash{}neg S) = 0.59\$.} \\
\texttt{4.  The probability of cholera given liking spicy food: \$P(C | S) = 0.33\$.} \\
\\ 
\texttt{Using the Law of Total Probability, we can calculate the overall probability of cholera contraction:} \\
\texttt{\$\$P(C) = P(C | S) \textbackslash{}cdot P(S) + P(C | \textbackslash{}neg S) \textbackslash{}cdot P(\textbackslash{}neg S)\$\$} \\
\\ 
\texttt{Substituting the given values:} \\
\texttt{\$\$P(C) = (0.33 \textbackslash{}cdot 0.41) + (0.59 \textbackslash{}cdot 0.59)\$\$} \\
\\ 
\texttt{Calculating the terms:} \\
\texttt{\$\$0.33 \textbackslash{}cdot 0.41 = 0.1353\$\$} \\
\texttt{\$\$0.59 \textbackslash{}cdot} \\
            }
        }
    }
    & & \\ \\

    \theutterance \stepcounter{utterance}  
    & & & \multicolumn{2}{p{0.3\linewidth}}{
        \cellcolor[rgb]{0.9,0.9,0.9}{
            \makecell[{{p{\linewidth}}}]{
                \texttt{\tiny{[GM$|$GM]}}
                \texttt{Yes} \\
            }
        }
    }
    & & \\ \\

    \theutterance \stepcounter{utterance}  
    & & & \multicolumn{2}{p{0.3\linewidth}}{
        \cellcolor[rgb]{0.9,0.9,0.9}{
            \makecell[{{p{\linewidth}}}]{
                \texttt{\tiny{[GM$|$GM]}}
                \texttt{no} \\
            }
        }
    }
    & & \\ \\

    \theutterance \stepcounter{utterance}  
    & & & \multicolumn{2}{p{0.3\linewidth}}{
        \cellcolor[rgb]{0.9,0.9,0.9}{
            \makecell[{{p{\linewidth}}}]{
                \texttt{\tiny{[GM$|$GM]}}
                \texttt{game\_result = LOSE} \\
            }
        }
    }
    & & \\ \\

\end{supertabular}
}

\end{document}
